\begin{theorem}[持续分支自然扩张的严格支持链与全深度 exact ledger 的可数无穷支撑]\label{thm:conclusion-primeslice-omega-support-chain-countable}
\leanverified{paper\_conclusion\_primeslice\_omega\_support\_chain\_countable}
设 \(T:X\twoheadrightarrow X\) 为可数集合上的有限纤维满射，并假设对每个 \(N\ge 1\) 的 \(N\)-步前史空间，其每一层都存在真实分支。记
\[
\mathfrak c_N
\]
为精确分离全部 \(N\)-步前历史在最坏情形下所需的最少非平凡 prime slices 数目。则
\[
\mathfrak c_N=N.
\]
从而可取一条严格递增的最优支持链
\[
\Sigma_1\subsetneq \Sigma_2\subsetneq \cdots,
\qquad
\#\Sigma_N=N,
\]
并且该链不存在有限 cofinal 子集。特别地，任意忠实的全深度 prime-register 账本实现都必须使用可数无穷多个彼此独立的 prime slices；尤其不存在任何有限支持的精确外置。
\end{theorem}
\begin{proof}
推论 \ref{cor:xi-time-part74-finite-past-primeslice-complexity-depth} 已给出
\[
\mathfrak c_N=N.
\]
对每个 \(N\)，取一个恰使用 \(N\) 个活跃 prime slices 的最优实现，并把其支持记为 \(\Sigma_N\)。由 one-prime-per-layer 的标准最优构造，可令这些支持满足
\[
\Sigma_1\subsetneq \Sigma_2\subsetneq \cdots,
\qquad
\#\Sigma_N=N.
\]
若存在有限 cofinal 子集 \(\Sigma\)，则 \(\#\Sigma<\infty\)。取 \(N>\#\Sigma\)，便不可能由 \(\Sigma\) 承载一个 \(N\)-步最优外置，这与 \(\mathfrak c_N=N\) 矛盾。故该链无有限 cofinal 子集。

最后，由定理 \ref{thm:conclusion-natural-extension-omega-layer-minimal-externalization}，持续分支系统的精确无限前史外置在最坏情形下必须使用 \(\omega\)-层 prime-slice 结构。结合上面的严格支持链，便得到任意忠实全深度账本都必须使用可数无穷支撑。证毕。
\end{proof}

\begin{proposition}[全深度 prime-support ledger 的实现层/结构层半圆维分裂]\label{prop:conclusion-primeslice-omega-support-impl-structure-split}
\leanverified{paper\_conclusion\_primeslice\_omega\_support\_impl\_structure\_split}
设
\[
\mathcal M_\omega:=\varinjlim_{N}\NN^{(\Sigma_N)}
\]
为定理 \ref{thm:conclusion-primeslice-omega-support-chain-countable} 中全深度最优支持链所诱导的直极限交换幺半群。则 \(\mathcal M_\omega\) 同构于可数生成自由交换幺半群，从而在选定一组素数标号后可与 \(\NN_{\times}\) 同构。特别地，
\[
\cdim^{\mathrm{impl}}(\mathcal M_\omega)=\frac12,
\qquad
\cdim_{+}^{\mathrm{hom}}(\mathcal M_\omega)=+\infty.
\]
因此，全深度 exact ledger 在集合实现层仍可压缩到单寄存器量级，而在同态保持层却不可有限线性化。
\end{proposition}
\begin{proof}
由定理 \ref{thm:conclusion-primeslice-omega-support-chain-countable}，支持链 \((\Sigma_N)\) 严格递增且其并为一个可数集。于是直极限
\[
\mathcal M_\omega=\varinjlim_{N}\NN^{(\Sigma_N)}
\]
正是这一可数生成元集合上的自由交换幺半群。选取该可数生成元集合与素数集 \(\mathbb P\) 的一个双射，便得到
\[
\mathcal M_\omega\cong \NN_{\times}
\]
的幺半群同构。

剩余两条便由命题 \ref{prop:killo-implementation-vs-structural-half-circle-dimension} 直接搬运而来：
\[
\cdim^{\mathrm{impl}}(\NN_{\times})=\frac12,
\qquad
\cdim_{+}^{\mathrm{hom}}(\NN_{\times})=+\infty.
\]
故对 \(\mathcal M_\omega\) 同样成立。证毕。
\end{proof}

\begin{theorem}[有限局部化支上的可见圆维恒一与连通 Lie 可见量对 slice growth 的盲性]\label{thm:conclusion-localization-visible-circle-blindness-slicegrowth}
\leanverified{paper\_conclusion\_localization\_visible\_circle\_blindness\_slicegrowth}
对任意有限活跃切片集 \(S\subset\mathbb P\)，记
\[
G_S:=\ZZ[S^{-1}].
\]
则 Pontryagin 对偶满足短正合列
\[
0\longrightarrow \prod_{p\in S}\ZZ_p
\longrightarrow \widehat{G_S}
\longrightarrow \TT
\longrightarrow 0,
\]
且
\[
\cdim_{\mathrm{fr}}(G_S)=1.
\]
因而当 \(S\subsetneq S'\) 时，可见连通圆因子始终保持为 \(\TT\)，而全部新增信息都严格落入全不连通核
\[
\prod_{p\in S'\setminus S}\ZZ_p
\]
之中。

进一步，若 \(\mathcal I\) 是任意只经由 \(\widehat{G_S}\) 的连通紧 Lie 商因子化的不变量，则
\[
\mathcal I(\widehat{G_S})
\]
实际上与 \(S\) 无关；亦即任何固定维的连通 Lie 可见量都对有限活跃 slice 的增长完全失明。
\end{theorem}
\begin{proof}
短正合列与
\[
\cdim_{\mathrm{fr}}(G_S)=1
\]
正是定理 \ref{thm:killo-localization-circle-dimension-prime-ledger-splitting} 的内容。因此，当 \(S\subsetneq S'\) 时，\(\widehat{G_S}\) 与 \(\widehat{G_{S'}}\) 的连通商都同构于 \(\TT\)，而新增素数支持只会扩大 profinite 核。

现设
\[
q:\widehat{G_S}\twoheadrightarrow L
\]
为任意连通紧 Lie 商。取 Pontryagin 对偶，得到单射
\[
L^\wedge\hookrightarrow G_S.
\]
由于 \(L\) 连通紧 Lie，\(L^\wedge\) 为有限生成无挠交换群，且
\[
\rank_{\ZZ}(L^\wedge)=\dim L.
\]
而
\[
\rank_{\ZZ}(G_S)=1,
\]
故必有
\[
\dim L\le 1.
\]
于是任意非平凡连通紧 Lie 商都只能是一维圆群 \(\TT\)；换言之，\(\widehat{G_S}\) 的全部连通 Lie 可见部分都塌缩到同一个圆商。故任何只经由连通 Lie 商因子化的不变量都不可能区分不同的有限 \(S\)。证毕。
\end{proof}

\endinput
